\documentclass[10pt,a4paper]{amsart}
\usepackage[margin=19mm]{geometry}
\usepackage{amsmath,amssymb,mathtools}
\usepackage[scr=rsfs,cal=euler]{mathalfa}
\usepackage{xfrac}
\usepackage[unicode,hidelinks,bookmarks=false]{hyperref}
\newcommand{\ie}{i.e.\ }
\newcommand{\cf}{cf.\ }
\newcommand{\resp}{respectively\ }
\newcommand{\Subsection}{\subsection}
\providecommand{\nobreakdash}{\nobreak\hbox{-}}
\setlength{\emergencystretch}{2em}
\multlinegap0pt
\newtheorem{theorem}{Theorem}
\newtheorem{lemma}{Lemma}
\newtheorem{assumption}{Assumption}
\title[Theorem 1]{An accelerated proximal PRS-SQP algorithm with dual ascent-descent procedures for smooth composite optimization (Theorem 1)}
\author{Jiachen Jin, Guodong Ma, Jinbao Jian}
\date{Source: arXiv:2505.09078v2 (CC BY 4.0)}
\begin{document}
\maketitle

\noindent\emph{Source and licence.} An accelerated proximal PRS-SQP algorithm with dual ascent-descent procedures for smooth composite optimization, by Jiachen Jin, Guodong Ma, Jinbao Jian. Version arXiv:2505.09078v2
(CC BY 4.0), distributed by arXiv under the Creative Commons Attribution 4.0 International licence,
http://creativecommons.org/licenses/by/4.0/, as recorded in the arXiv licence metadata for this exact
version and shown on the abstract page https://arxiv.org/abs/2505.09078v2. Full licence notice:
\texttt{licenses/arxiv-2505.09078-CC-BY-4.0.txt}.\par\smallskip

\noindent\textit{Editorial scope.} Counted unit: the article's Theorem 1 (source0.tex lines 724--735): for the composite problem under the article's assumptions, every accumulation point of the sequence generated by the accelerated proximal PRS-SQP algorithm with dual ascent-descent procedures is a Karush-Kuhn-Tucker point. The article's complete original proof of it is delivered with it (source0.tex lines 736--777, one \texttt{proof} environment of 42 lines, adjacent to the statement, with no gap and no deferral). No mathematics is written, repaired or re-derived: the statement and the proof are the article's own lines, in the article's own notation, and no retained line is substituted, re-flowed or omitted, so nothing is removed and the package carries no omission record. Retained uncounted as the source-local context the proof needs: the six displayed defining relations of the algorithm and of its optimality conditions (lines 311--313, 319--321, 331--333, 342--344, 350--352 and 364--366), the first and third standing assumptions (lines 430--441 and 644--646, source labels \texttt{ass1} and \texttt{ass3}), the three lemmas bounding the sequence, the multipliers and the proximity terms (lines 503--505, 529--536 and 561--565, source labels \texttt{lemNA}, \texttt{lem2} and \texttt{lem3}) and the displayed KKT-type relation they produce (lines 608--614, source label \texttt{LK1LK}). Every cross-reference of every retained line resolves inside this document, and no retained line contains a citation, so the wrapper carries no bibliography and no bibliography entry. No retained line contains a figure environment or an image-inclusion command; the article's own two figures lie outside every retained span, so no image is delivered and the package declares no asset dependency.

Publisher class: the article is typeset with the Springer Nature class \texttt{sn-jnl}, whose class file \texttt{sn-jnl.cls} is shipped in the e-print. That class file is not redistributed in this package and is not used by it: the wrapper is an AMS \texttt{amsart} document that restates the article's own three environment declarations, namely \texttt{theorem}, \texttt{lemma} and \texttt{assumption}, each on its own counter exactly as the article declares them, together with the standard packages the retained lines need. The retained environments therefore print in this document as Assumption 1, Lemma 1, Lemma 2, Lemma 3, Assumption 2 and Theorem 1 in order of appearance, whereas the article prints the same six items with the numbering of its own full document; the numbered displays of the retained lines print consecutively in order of appearance, whereas the article numbers its equations per section.

The complete native source of the article as distributed in the arXiv e-print for this exact version is bundled unchanged as \texttt{sources/178090/source0.tex}: it is byte identical to the e-print file \texttt{HAP-PRS-SQP.tex} except that the pipeline's mechanical concatenation prepends one header comment line naming it, and except that the single non-UTF-8 byte of that file is carried through the pipeline's lossy decode exactly as the pipeline read it.
\begin{equation}\label{xoc}
\nabla f(x_{k})-A^{\top}(\lambda_{k}-\beta(A x_{k}-y_{k}))+\mathcal{H}_k^x(\tilde{x}_{k+1}-x_{k})=0.
\end{equation}

\begin{equation}\label{sd-x}
d_{k}^{x} = \bar{x}_{k+1}-x_{k}=(1+\alpha)(\tilde{x}_{k+1}-x_{k}).
\end{equation}

\begin{equation}\label{lambda1}
\lambda_{k+\frac{1}{2}} = \lambda_{k}-r \beta(A x_{k+1}-y_{k}).
\end{equation}

\begin{equation}\label{yoc}
\nabla g (y_{k})+\lambda_{k+\frac{1}{2}}-\beta(A x_{k+1}-y_{k})+\mathcal{H}_k^y(\tilde{y}_{k+1}-y_{k})=0.
\end{equation}

\begin{equation}\label{sd-y}
d_{k}^{y} = \bar{y}_{k+1}-y_{k}=(1+\alpha)(\tilde{y}_{k+1}-y_{k}).
\end{equation}

\begin{equation}\label{lambda2}
\lambda_{k+1} = \lambda_{k+\frac{1}{2}}-s \beta(A x_{k+1}-y_{k+1}),
\end{equation}

\begin{assumption}\label{ass1}{\rm
		
		(i) The iteration sequence $\{w_k\}$ and matrix sequences $\{{H}_k^x\}$ and $\{{H}_k^y\}$ generated by the HAP-PRS-SQP algorithm are bounded, with $\|{H}_k^x\|\leq \eta^x$ and $\|{H}_k^y\|\leq \eta^y$ for all $k$.
		
		(ii) There exist lower bounds $\underline{\lambda}^x$ and $\underline{\lambda}^y$ of $\{\lambda_{\min}(H^x_k)\}$ and $\{\lambda_{\min}(H^y_k)\}$, respectively, such that
		\begin{equation}\label{NF01} \eta^x_1:=\underline{\lambda}^x+\beta\lambda_{\min}(A^TA)+\ell>0,\
			\eta^y_1:=\underline{\lambda}^y+\beta+\sigma>0.
		\end{equation}
		
		(iii) The gradients $\nabla f(x)$ and $\nabla g(y)$ are local Lipschitz continuous.
	}
\end{assumption}

\begin{lemma}\label{lemNA}
	If Assumption \ref{ass1} (i,ii) holds, then the sequences $\{d_k\}$, $\{(\tilde{x}_k,\tilde{y}_k)\}$ and $\{(\bar{x}_k,\bar{y}_k)\}$ generated by HAP-PRS-SQP are all bounded, and so is the sequence $\{\hat{w}_k\}$.
\end{lemma}

\begin{lemma}\label{lem2}
	If Assumption \ref{ass1}  holds, then the sequence $\{t_k\}$ yielded by the HAP-PRS-SQP algorithm possess a positive infimum independent of $k$ as follows:
	\[\label{stepL}
	\min \{t_k^x,t_k^y\} \geq \nu\min \left\{1,\frac{(\frac{1}{1+\alpha}-\rho)\eta_1^{x}}
	{L_f+\beta\|A^{\top}A\|},\frac{(\frac{1}{1+\alpha}-\rho)\eta_1 ^y}{L_g+\beta}\right\}:=\gamma,
	\]
	where $\rho,\nu\in (0,1)$ and $\alpha \in (-1, \frac{1}{\rho}-1)$.
\end{lemma}

\begin{lemma}\label{lem3}
	Suppose that Assumption \ref{ass1} holds. Then
	\[\hat{L}_{\beta}(\hat{w}_{k+1})-\hat{L}_{\beta}(\hat{w}_{k}) \leq -\delta^x\|d_{k}^x\|^2-\delta^y\|d_{k}^y\|^2,\]
	where $\delta^x=\rho\gamma\eta^x_1-\frac{6(1-s)^2\beta\lambda_{\max}(A^{\top}A)}{|r+s|}$, $\delta^y =\rho\gamma\eta^y_1 -\frac{6}{|r+s|\beta}\left(L_{g}^{2}+(1+s^{2})\beta^{2}+2\left(\frac{\eta_2^y}{1+\alpha}\right)^2\right)-\frac{|rs|\beta}{|r+s|}$.
\end{lemma}

	\begin{align}
		\frac{1}{6}\|\lambda_{k+1}-\lambda_{k}\|^{2}
		&\leq(L_{g}^{2}+\beta^2)\|y_{k}-y_{k-1}\|^{2}+(1-s)^{2}\beta^{2}\|A(x_{k+1}-x_{k})\|^{2}\nonumber\\
		&\ \ \ \ +s^{2}\beta^{2}\|y_{k+1}-y_{k}\|^{2}  +\left(\frac{\eta_2^y}{1+\alpha}\right)^2\|d_k^y\|^2+\left(\frac{\eta_2^y}{1+\alpha}\right)^2\|d_{k-1}^y\|^2 \nonumber\\
		&\leq(L_{g}^{2}+\beta^2)\|d_{k-1}^y\|^{2}+(1-s)^{2}\beta^{2}\lambda_{\max}(A^{\top}A)\|d_{k}^x\|^{2}\nonumber\\
		&\ \ \ \ + s^{2}\beta^{2}\|d_{k}^y\|^{2}+\left(\frac{\eta_2^y}{1+\alpha}\right)^2\|d_k^y\|^2+\left(\frac{\eta_2^y}{1+\alpha}\right)^2\|d_{k-1}^y\|^2. \label{LK1LK}
	\end{align}

\begin{assumption}\label{ass3}{\rm Suppose that the parameters in the HAP-PRS-SQP algorithm are chosen such that
$\delta^x> 0$ and $\delta^y> 0$. 
}\end{assumption}

\begin{theorem}\label{the1}(Global convergence) Let $\Omega$ and $\hat{\Omega}$ denote the cluster point set of sequences $\{w_{k}\}$ and $\{\hat{w}_{k}\}$, respectively. Suppose that Assumptions \ref{ass1} and \ref{ass3} hold. Then the five claims for the HAP-PRS-SQP algorithm as follows hold true.

(i) The whole sequence $\{\hat{L}_{\beta}(\hat{w}_k)\}$ is convergent, and $\hat{L}_{\beta}(\hat{w}_*)=\lim_{k \to \infty }\hat{L}_{\beta}(\hat{w}_k)=\inf_{k}\hat{L}_{\beta}(\hat{w}_k)$, $\forall~\hat{w}_* \in \hat{\Omega}$. Therefore, $\hat{L}_{\beta}(\cdot)$ is finite and constant on $\hat{\Omega}$.

(ii) $ \lim_{k\to\infty}\|d_k^x\|=\lim_{k\to\infty}\|d_k^y\|=0,~\sum_{k=0}^{+\infty}\| w_{k+1}-w_k \|^2<+\infty$.

(iii) $\Omega$ and $\hat{\Omega}$ are nonempty compact sets, and  $d(w_k, \Omega) \to 0$ and $d(\hat{w}_k, \hat{\Omega}) \to 0$ as $k \to \infty$.

(iv) $\Omega \subseteq {\rm crit}~L_{\beta}$, so the HAP-PRS-SQP algorithm is global convergent.

(v) $\hat{\Omega} = \{(x_*,y_*,\lambda_*,0):~(x_*,y_*,\lambda_*) \in \Omega\}=\Omega\times {0}$.
\end{theorem}

\begin{proof}
(i) By Lemma \ref{lemNA}, the sequence $\{\hat{w}_{k}\}$ is bounded, and so is $\{\hat{L}_{\beta}(\hat{w}_{k})\}$. Combined with the monotonic descent of $\{\hat{L}_{\beta}(\hat{w}_{k})\}$ (Lemma \ref{lem3}), it follows that $\{\hat{L}_{\beta}(\hat{w}_{k})\}$ is convergent. For any $\hat{w}^*\in \hat{\Omega}$, there exists a subsequence $\{\hat{w}_{k_j}\}$ such that $\hat{w}_{k_j} \to \hat{w}_*$, $k_j\to\infty$. Thus,
$\hat{L}_{\beta}(\hat{w}_*)=\lim_{k_j \to \infty }\hat{L}_{\beta}(\hat{w}_{k_j})=\lim_{k \to \infty }\hat{L}_{\beta}(\hat{w}_{k})=\inf_{k}\hat{L}_{\beta}(\hat{w}_k)$.

(ii) From Lemma \ref{lem3}, we have
\[	\delta^x\|d_{k}^x\|^2+\delta^y\|d_{k}^y\|^2\leq\hat{L}_{\beta}(\hat{w}_{k}) -\hat{L}_{\beta}(\hat{w}_{k+1}).\]
Summing this inequality for $k=0,1, \dots, p$ and noticing that $\hat{L}_{\beta}(\hat {w}_{0})<+\infty$, we obtain
\[	
\sum_{k=0}^p(\delta^x\|d_{k}^x\|^2+\delta^y\|d_{k}^y\|^2)\leq\hat{L}_{\beta}(\hat{w}_{0}) -\hat{L}_{\beta}(\hat{w}_{p+1})
\leq\hat{L}_{\beta}(\hat{w}_{0}) -\hat{L}_{\beta}(\hat{w}_{*})<+\infty.
\]
Since $\delta^x>0$ and $\delta^y>0$, it follows that
\[\sum_{k=0}^{+\infty} \|d_k^x\|^2<+\infty,~\sum_{k=0}^{+\infty} \|d_k^y\|^2<+\infty.\]
Using \eqref{LK1LK}, we further deduce
\[
\sum_{k=0}^{+\infty} \|{\lambda}_{k+1}-\lambda_{k}\|^2<+\infty;\ \lim_{k\to\infty}\|d_k^x\|=\lim_{k\to\infty}\|d_k^y\|=0.
\]
Considering that
\begin{equation}\label{xleqd}
\|x_{k+1}-x_k\|\leq \frac{1}{t_k^x}\|{x}_{k+1}-x_{k}\|=\|d_k^x\|,\ \|y_{k+1}-y_k\|\leq \frac{1}{t_k^y}\|{y}_{k+1}-y_{k}\|=\|d_k^y\|.
\end{equation}
Thus,
\[\sum_{k=0}^{+\infty}\|x_{k+1}-x_k\|^2<+\infty,~\sum_{k=0}^{+\infty}\|y_{k+1}-y_k\|^2<+\infty.\]
Consequently, $\sum_{k=0}^{+\infty}\|w_{k+1}-w_k\|^2<+\infty$.

(iii) The boundedness of $\{w_{k}\}$ and $\{\hat{w}_{k}\}$, along with the definitions of $\Omega$ and $\hat{\Omega}$, implies the claim directly.

(iv) For any given $w_*=(x_*,y_*,\lambda_*)\in\Omega$, there exists a subsequence $\{w_{k_j}\}$ such that $w_{k_j} \to w_*$, $k_j \to \infty$. Since $\|w_{k+1}-w_k\|\to 0$ as $k\to\infty$, it follows that $w_{k_j+1} \to w_*$, $k_j \to \infty$. From \eqref{lambda1}, $\{\lambda_{{k_j}+\frac{1}{2}}\}$ is convergent.
Let $\lambda_{{k_j}+\frac{1}{2}} \to \lambda_{**}$ as $k_j\to \infty$. Taking the limit $k={k_j}\to \infty$ in \eqref{lambda1} and \eqref{lambda2}, we obtain
\[\lambda_{**}=\lambda_{*}-r\beta(Ax_*-y_*),~\lambda_{*}=\lambda_{**}-s\beta(Ax_*-y_*).\]
Since $r+s\not=0$ and $\beta>0$, it follows that $Ax_*=y_*$ and $\lambda_{**}=\lambda_{*}$.
From \eqref{sd-x}, \eqref{sd-y} and conclusion (ii), we have
\[\lim\limits_{k_j\to \infty}\tilde{x}_{{k_j}+1}=\lim\limits_{k_j\to \infty}(x_{k_j}+\frac{1}{1+\alpha}d_{k_j}^x)=x_*,
\lim\limits_{k_j\to \infty}\tilde{y}_{{k_j}+1}=\lim\limits_{k_j\to \infty}(y_{k_j}+\frac{1}{1+\alpha}d_{k_j}^y)=y_*.\]
By the continuity of $\nabla f$ and $\nabla g$, and using $Ax_*=y_*$, taking the limit in \eqref{xoc} and \eqref{yoc} as $k=k_j\to \infty$, we obtain
\[
\nabla f(x_*)-A^{\top}\lambda_*=0,~\nabla g(y_*)+\lambda_*=0,~Ax_*-y_*=0.
\]
Thus, $w_*=(x_*,y_*,\lambda_*)$ satisfies the KKT conditions and belongs to crit $L_\beta$.

(v) By the definition of $\hat{w}_k$ and $d_k^y \to 0$ as $k \to \infty$, the claim follows directly.
\end{proof}

\end{document}
